\section{Introduction}

Ce laboratoire porte sur l'identification expérimentale des paramètres d'un actionneur à bobine mobile. Ce type d'actionneur est un transducteur électromécanique à déflexion angulaire limitée, largement utilisé en instrumentation de précision : galvanomètres, têtes d'impression, scanners laser ou actionneurs de disques durs.

Le système est entièrement caractérisé par six paramètres physiques :
\begin{itemize}
    \item \(R\) : résistance de la bobine [\si{\ohm}] ;
    \item \(L\) : inductance de la bobine [\si{H}] ;
    \item \(K_u = K_t\) : constante électromécanique de couplage [\si{V.s/rad} = \si{N.m/A}] ;
    \item \(J\) : moment d'inertie de la partie mobile [\si{kg.m^2}] ;
    \item \(C_v\) : coefficient d'amortissement visqueux [\si{N.m.s/rad}] ;
    \item \(K_r\) : raideur du rappel élastique [\si{N.m/rad}].
\end{itemize}

Trois expériences complémentaires sont menées pour identifier ces paramètres :
\begin{itemize}
    \item \textbf{EX1} — réponse indicielle avec et sans mouvement : identification de \(R\) et \(L\) ;
    \item \textbf{EX2} — oscillations libres bobine ouverte : identification de \(C_v/J\), \(K_r/J\) et \(K_u\) ;
    \item \textbf{EX3} — oscillations libres bobine court-circuitée : identification de \(J\), \(C_v\) et \(K_r\).
\end{itemize}

Les signaux sont acquis à l'oscilloscope Tektronix DPO2014B et traités sous Python dans le notebook \texttt{traitement.ipynb}.

\paragraph{Données de l'actionneur.}
L'angle de déflexion maximal est \(\Delta\alpha_{\max} = \SI{17.1}{\degree} = \SI{0.298}{rad}\).
La sonde de courant utilisée a un gain de \(\SI{1}{A/V}\).

\paragraph{Appareillage utilisé.}
\begin{table}[H]
\centering
\caption{Instruments utilisés}
\begin{tabular}{ll}
\toprule
Appareil & Modèle \\
\midrule
Oscilloscope        & Tektronix DPO2014B \\
Sonde de courant    & gain \(\SI{1}{A/V}\) \\
Alimentation DC     & GW INSTEK PPE-3323 \\
Actionneur          & bobine mobile à déflexion limitée \\
\bottomrule
\end{tabular}
\label{tab:appareils}
\end{table}
